\documentclass[10pt]{article}
\usepackage[T1]{fontenc}
\usepackage[utf8]{inputenc}
\usepackage[french]{babel}
\usepackage{amsmath,amssymb,mathrsfs,mathpazo}
\usepackage{geometry}
\geometry{a4paper,margin=2cm}
\usepackage{graphicx,array,tabularx,tabularray,colortbl,multirow,multicol,booktabs,enumitem,pifont,ulem,fancybox,diagbox,cancel}
\usepackage{pstricks,pstricks-add,pst-plot,pst-tree,pst-node,pst-text,pst-func,pst-math,pst-bspline,pst-3dplot,pst-eucl,pst-solides3d}
\let\pscalculate\undefined
\usepackage{pst-bezier}
\usepackage{tikz,pgfplots,tkz-tab}
\usetikzlibrary{arrows,arrows.meta,patterns,calc,positioning,trees,decorations.markings,shapes.geometric,shapes.misc,shapes.arrows,intersections,decorations.pathmorphing,decorations.pathreplacing,backgrounds,shadows,babel}
\pgfplotsset{compat=1.18}
\usepgfplotslibrary{fillbetween}
\usepackage[np]{numprint}
\usepackage{listings,makecell,pas-tableur,twoopt}
\let\barre\undefined
\usepackage[pstricks]{tablvar}
\definecolor{darkgreen}{RGB}{0,100,0}
\providecommand{\red}{\color{red}}
\providecommand{\blue}{\color{blue}}
\providecommand{\green}{\color{green}}
\providecommand{\black}{\color{black}}
\providecommand{\white}{\color{white}}
\providecommand{\gray}{\color{gray}}
\providecommand{\cyan}{\color{cyan}}
\providecommand{\magenta}{\color{magenta}}
\providecommand{\yellow}{\color{yellow}}
\providecommand{\violet}{\color{violet}}
\providecommand{\orange}{\color{orange}}
\providecommand{\euro}{\texteuro}
\providecommand{\EUR}{\texteuro}
\providecommand{\ieme}{\textsuperscript{e}}
\providecommand{\degres}{\ensuremath{^\circ}}
\providecommand{\degree}{\ensuremath{^\circ}}
\providecommand{\vv}[1]{\overrightarrow{#1}}
\providecommand{\Coord}[1]{#1}
\providecommand{\decosix}{}
\providecommand{\starredbullet}{\ensuremath{\bullet}}
\pagestyle{empty}
\setlength{\parindent}{0pt}
\usepackage{adjustbox}\DeclareUnicodeCharacter{2236}{:}\DeclareUnicodeCharacter{2190}{\ensuremath{\leftarrow}}
\providecommand{\R}{\mathbb{R}}
\providecommand{\N}{\mathbb{N}}
\providecommand{\Z}{\mathbb{Z}}
\providecommand{\Q}{\mathbb{Q}}
\providecommand{\D}{\mathbb{D}}
\providecommand{\C}{\mathbb{C}}
\providecommand{\e}{{\mathrm{e}}}
\providecommand{\vect}[1]{\overrightarrow{#1}}
\providecommand{\vectt}[1]{\overrightarrow{\mathrm{#1}}}
\providecommand{\vvt}[1]{\overrightarrow{\mathrm{#1}}}
\providecommand{\np}[2][]{{#2~#1}}
\providecommand{\vv}[1]{{\overrightarrow{#1}}}
\renewcommand{\d}{{\mathrm{d}}}
\providecommand{\degres}{{^\circ}}
\providecommand{\degre}{{^\circ}}
\providecommand{\blue}{\color{blue}}
\providecommand{\red}{\color{red}}
\providecommand{\enskip}{\quad}
\providecommand{\llbracket}{\lbrack\!\lbrack}
\providecommand{\rrbracket}{\rbrack\!\rbrack}
\providecommand{\euro}{{\text{€}}}
\providecommand{\no}{{\text{n°}}}
\providecommand{\interval}[2]{{[#1;#2]}}
\providecommand{\intervFO}[2]{{[#1;#2[}}
\providecommand{\dotfill}{\ldots}

\providecommand{\ds}{\displaystyle}
\providecommand{\textding}[1]{#1}
\providecommand{\Oij}{\(\mathrm{O};\vec{\imath},\vec{\jmath}\)}
\providecommand{\Oijk}{\(\mathrm{O};\vec{\imath},\vec{\jmath},\vec{k}\)}
\providecommand{\Ouv}{\(\mathrm{O};\vec{u},\vec{v}\)}
\providecommand{\intervFF}[2]{\ensuremath{[#1;#2]}}
\renewcommand{\euro}{\ifmmode\text{\texteuro}\else\texteuro\fi}
\DeclareUnicodeCharacter{00B0}{\ifmmode^{\circ}\else\textdegree\fi}
\DeclareUnicodeCharacter{20AC}{\ifmmode\text{\texteuro}\else\texteuro\fi}
% Énoncés : Éducation nationale. Transcription et corrigés : APMEP, Denis Vergès et contributeurs. Découpage et classement : CoopMaths, Sébastien Lozano, Éric Elter et contributeurs. Adaptation de la présentation : Maths Explorer.
% CC BY-SA : https://coopmaths.fr/www/about/
% https://coopmaths.fr/alea/static/e3c/2020/tex/e3c_2020_00_sujet52_2.tex
\begin{document}
\noindent\textit{APMEP / CoopMaths. CC BY-SA. Présentation adaptée par Maths Explorer.}\par\medskip
\medskip

Une fleuriste met en vente quatre sortes de bouquets dont les tarifs et la composition sont
indiqués dans le tableau ci-dessous :

\medskip
\begin{center}
\begin{tabularx}{0.85\linewidth}{|*{2}{X|}}\hline
Bouquet de tulipes orange : 10,50 \euro  & Bouquet de roses orange : 23,50 \euro\\\hline
Bouquet de tulipes blanches : 11,60 \euro& Bouquet de roses blanches : 25,50 \euro\\\hline
\end{tabularx}
\end{center}
\medskip

\begin{itemize}
\item 72\,\% des bouquets mis en vente ne contiennent que des roses.
\item Les autres bouquets mis en vente ne contiennent que des tulipes.
\item 20\,\% des bouquets de tulipe mis en vente ne contiennent que des tulipes orange.
\item 36\,\% des bouquets mis en vente ne contiennent que des roses blanches.
\end{itemize}
Un client achète au hasard un bouquet parmi ceux mis en vente par la fleuriste. On note :

\begin{itemize}
\item $R$ l'évènement : \og Le bouquet acheté par ce client est composé de roses.\fg
\item $B$ l'évènement : \og Le bouquet acheté par ce client est composé de fleurs blanches. \fg
\end{itemize}
Les évènements contraires des évènements $R$ et $B$ sont notés respectivement $\overline{R}$ et $\overline{B}$.

\medskip

\begin{enumerate}
\item 
	\begin{enumerate}
		\item Donner, sans justifier, la probabilité $p(R \cap B)$.
		\item Recopier et compléter le plus possible l'arbre de probabilité ci-dessous en
traduisant uniquement les données de l'énoncé.
\begin{center}

\psset{nodesepA=0pt,nodesepB=3pt,treesep=0.75,labelsep=0.1pt,levelsep=2.5cm}
\pstree[treemode=R]{\TR{}}
{\pstree{\TR{}\taput{}}
	{
	\TR{}\taput{}
	\TR{}\tbput{}
	}
\pstree{\TR{~}\tbput{}}
	{\TR{}\taput{}
	\TR{}\tbput{}
	}
}
\end{center}

		\item Montrer que $p(B) = 0,584$.
	\end{enumerate}
\item On note $X$ la variable aléatoire qui donne le prix d'un bouquet acheté par un client.
	\begin{enumerate}
		\item  Recopier et compléter le tableau ci-dessous donnant, pour chaque valeur $x_i$ de $X$,
la probabilité de l'évènement $\{X = x_i \}$. Justifier.

\begin{center}
\begin{tabularx}{0.6\linewidth}{|m{1.82cm}|*{4}{>{\centering \arraybackslash}X|}}\hline
\centering $x_i$		&	&	&	& \\\hline
\centering $p(X = x_i)$ &	&	&	&\\\hline
\end{tabularx}
\end{center}

		\item Calculer l'espérance de la variable aléatoire $X$. \emph{On arrondira le résultat au centième.}
	\end{enumerate}
\end{enumerate}

\vspace{0,5cm}

\end{document}
